% mcmicm_styles.tex
% 从美赛模板提取的代码/AI 报告样式
% 使用时导言区加入: % mcmicm_styles.tex
% 从美赛模板提取的代码/AI 报告样式
% 使用时导言区加入: % mcmicm_styles.tex
% 从美赛模板提取的代码/AI 报告样式
% 使用时导言区加入: % mcmicm_styles.tex
% 从美赛模板提取的代码/AI 报告样式
% 使用时导言区加入: \input{mcmicm_styles.tex}

\usepackage{tcolorbox} % 美丽盒子包
\usepackage{matlab-prettifier} % MATLAB代码排版包
\tcbuselibrary{listings,skins,breakable,xparse}

% 颜色定义
\definecolor{darkgray}{rgb}{0.4,0.4,0.4}
\definecolor{darkblue}{rgb}{0.0,0.0,0.5}
\definecolor{darkgreen}{rgb}{0,0.45,0.13}
\definecolor{orange}{rgb}{0.87,0.49,0.0}

% ========== tcolorbox 列表样式 ==========
\tcbset{
    mylist/.style = {
        colframe = gray,
        colback = white,
        coltitle = red!50!yellow!3!white,
        colbacktitle = white,
        listing only,
        attach boxed title to top center = {yshift = -\tcboxedtitleheight/2},
        enhanced,
        drop fuzzy shadow,
        left = 6.5mm,
        breakable,
        fonttitle=\small\bfseries\color{black},
        before skip=12pt plus 2pt,
        after skip=12pt plus 2pt,
    },
    example/.style 2 args = {
        mylist,
        title = {Listing \thetcbcounter: #1},
        label = {#2},
    },
}

\newtcblisting[auto counter, number within = section, list inside = mcode]{matlab}[3][]{%
    listing options = {
        style = Matlab-editor,
        numbers = left,
        numberstyle = \footnotesize\color{darkgray}\ttfamily,
        basicstyle = \footnotesize\ttfamily,
        breaklines = true,
        breakatwhitespace = true,
        postbreak = \mbox{\textcolor{gray}{$\hookrightarrow$}\space}
    },
    overlay = {\begin{tcbclipinterior}\fill[blue!15!white] (frame.south west) rectangle ([xshift=5.3mm]frame.north west);\end{tcbclipinterior}},
    example = {#2}{#3}, #1,
}

\newtcblisting[auto counter, number within = section, list inside = pcode]{python}[3][]{%
    listing options = {
        language = Python,
        numbers = left,
        numberstyle = \footnotesize\color{darkgray}\ttfamily,
        basicstyle = \footnotesize\ttfamily,
        keywordstyle = \color{darkblue}\bfseries,
        commentstyle = \color{darkgreen},
        stringstyle = \color{red!60},
        showstringspaces = false,
        emphstyle = \color{orange!60},
        emph = {self, cls, True, False, None},
        breaklines = true,
        breakatwhitespace = true,
        postbreak = \mbox{\textcolor{gray}{$\hookrightarrow$}\space}
    },
    overlay = {\begin{tcbclipinterior}\fill[blue!15!white] (frame.south west) rectangle ([xshift=5.3mm]frame.north west);\end{tcbclipinterior}},
    example = {#2}{#3}, #1,
}

% ========== AI 报告专用无边框 listings 风格 ==========
\lstdefinestyle{ai_python_style}{
    language = Python,
    numbers = none,
    basicstyle = \footnotesize\ttfamily,
    keywordstyle = \color{darkblue}\bfseries,
    commentstyle = \color{darkgreen},
    stringstyle = \color{red!60},
    showstringspaces = false,
    emphstyle = \color{orange!60},
    emph = {self, cls, True, False, None},
    breaklines = true,
    breakatwhitespace = true,
    postbreak = \mbox{\textcolor{gray}{$\hookrightarrow$}\space},
    frame = none,
    backgroundcolor = \color{white},
    xleftmargin = 0pt,
    xrightmargin = 0pt,
}

\lstdefinestyle{ai_matlab_style}{
    style = Matlab-editor,
    numbers = none,
    basicstyle = \footnotesize\ttfamily,
    breaklines = true,
    breakatwhitespace = true,
    postbreak = \mbox{\textcolor{gray}{$\hookrightarrow$}\space},
    frame = none,
    backgroundcolor = \color{white},
    xleftmargin = 0pt,
    xrightmargin = 0pt,
}

\lstnewenvironment{ai_python}[1][]
    {\lstset{style=ai_python_style, #1}}
    {}

\lstnewenvironment{ai_matlab}[1][]
    {\lstset{style=ai_matlab_style, #1}}
    {}

% ========== 为定理类环境添加底色框 ==========
\mdfdefinestyle{thmshaded}{%
    backgroundcolor=gray!8,% 浅灰底色
    linecolor=black!30,% 边框颜色
    leftline=true,% 左侧线条
    topline=false,
    rightline=false,
    bottomline=false,
    linewidth=3pt,% 左侧线宽
    innerleftmargin=10pt,
    innertopmargin=6pt,
    innerbottommargin=6pt,
    roundcorner=0pt,
}

\surroundwithmdframed[style=thmshaded]{definition}
\surroundwithmdframed[style=thmshaded]{theorem}
\surroundwithmdframed[style=thmshaded]{lemma}
\surroundwithmdframed[style=thmshaded]{corollary}
\surroundwithmdframed[style=thmshaded]{assumption}
\surroundwithmdframed[style=thmshaded]{example}
\surroundwithmdframed[style=thmshaded]{problem}
\surroundwithmdframed[style=thmshaded]{solution}




\usepackage{tcolorbox} % 美丽盒子包
\usepackage{matlab-prettifier} % MATLAB代码排版包
\tcbuselibrary{listings,skins,breakable,xparse}

% 颜色定义
\definecolor{darkgray}{rgb}{0.4,0.4,0.4}
\definecolor{darkblue}{rgb}{0.0,0.0,0.5}
\definecolor{darkgreen}{rgb}{0,0.45,0.13}
\definecolor{orange}{rgb}{0.87,0.49,0.0}

% ========== tcolorbox 列表样式 ==========
\tcbset{
    mylist/.style = {
        colframe = gray,
        colback = white,
        coltitle = red!50!yellow!3!white,
        colbacktitle = white,
        listing only,
        attach boxed title to top center = {yshift = -\tcboxedtitleheight/2},
        enhanced,
        drop fuzzy shadow,
        left = 6.5mm,
        breakable,
        fonttitle=\small\bfseries\color{black},
        before skip=12pt plus 2pt,
        after skip=12pt plus 2pt,
    },
    example/.style 2 args = {
        mylist,
        title = {Listing \thetcbcounter: #1},
        label = {#2},
    },
}

\newtcblisting[auto counter, number within = section, list inside = mcode]{matlab}[3][]{%
    listing options = {
        style = Matlab-editor,
        numbers = left,
        numberstyle = \footnotesize\color{darkgray}\ttfamily,
        basicstyle = \footnotesize\ttfamily,
        breaklines = true,
        breakatwhitespace = true,
        postbreak = \mbox{\textcolor{gray}{$\hookrightarrow$}\space}
    },
    overlay = {\begin{tcbclipinterior}\fill[blue!15!white] (frame.south west) rectangle ([xshift=5.3mm]frame.north west);\end{tcbclipinterior}},
    example = {#2}{#3}, #1,
}

\newtcblisting[auto counter, number within = section, list inside = pcode]{python}[3][]{%
    listing options = {
        language = Python,
        numbers = left,
        numberstyle = \footnotesize\color{darkgray}\ttfamily,
        basicstyle = \footnotesize\ttfamily,
        keywordstyle = \color{darkblue}\bfseries,
        commentstyle = \color{darkgreen},
        stringstyle = \color{red!60},
        showstringspaces = false,
        emphstyle = \color{orange!60},
        emph = {self, cls, True, False, None},
        breaklines = true,
        breakatwhitespace = true,
        postbreak = \mbox{\textcolor{gray}{$\hookrightarrow$}\space}
    },
    overlay = {\begin{tcbclipinterior}\fill[blue!15!white] (frame.south west) rectangle ([xshift=5.3mm]frame.north west);\end{tcbclipinterior}},
    example = {#2}{#3}, #1,
}

% ========== AI 报告专用无边框 listings 风格 ==========
\lstdefinestyle{ai_python_style}{
    language = Python,
    numbers = none,
    basicstyle = \footnotesize\ttfamily,
    keywordstyle = \color{darkblue}\bfseries,
    commentstyle = \color{darkgreen},
    stringstyle = \color{red!60},
    showstringspaces = false,
    emphstyle = \color{orange!60},
    emph = {self, cls, True, False, None},
    breaklines = true,
    breakatwhitespace = true,
    postbreak = \mbox{\textcolor{gray}{$\hookrightarrow$}\space},
    frame = none,
    backgroundcolor = \color{white},
    xleftmargin = 0pt,
    xrightmargin = 0pt,
}

\lstdefinestyle{ai_matlab_style}{
    style = Matlab-editor,
    numbers = none,
    basicstyle = \footnotesize\ttfamily,
    breaklines = true,
    breakatwhitespace = true,
    postbreak = \mbox{\textcolor{gray}{$\hookrightarrow$}\space},
    frame = none,
    backgroundcolor = \color{white},
    xleftmargin = 0pt,
    xrightmargin = 0pt,
}

\lstnewenvironment{ai_python}[1][]
    {\lstset{style=ai_python_style, #1}}
    {}

\lstnewenvironment{ai_matlab}[1][]
    {\lstset{style=ai_matlab_style, #1}}
    {}

% ========== 为定理类环境添加底色框 ==========
\mdfdefinestyle{thmshaded}{%
    backgroundcolor=gray!8,% 浅灰底色
    linecolor=black!30,% 边框颜色
    leftline=true,% 左侧线条
    topline=false,
    rightline=false,
    bottomline=false,
    linewidth=3pt,% 左侧线宽
    innerleftmargin=10pt,
    innertopmargin=6pt,
    innerbottommargin=6pt,
    roundcorner=0pt,
}

\surroundwithmdframed[style=thmshaded]{definition}
\surroundwithmdframed[style=thmshaded]{theorem}
\surroundwithmdframed[style=thmshaded]{lemma}
\surroundwithmdframed[style=thmshaded]{corollary}
\surroundwithmdframed[style=thmshaded]{assumption}
\surroundwithmdframed[style=thmshaded]{example}
\surroundwithmdframed[style=thmshaded]{problem}
\surroundwithmdframed[style=thmshaded]{solution}




\usepackage{tcolorbox} % 美丽盒子包
\usepackage{matlab-prettifier} % MATLAB代码排版包
\tcbuselibrary{listings,skins,breakable,xparse}

% 颜色定义
\definecolor{darkgray}{rgb}{0.4,0.4,0.4}
\definecolor{darkblue}{rgb}{0.0,0.0,0.5}
\definecolor{darkgreen}{rgb}{0,0.45,0.13}
\definecolor{orange}{rgb}{0.87,0.49,0.0}

% ========== tcolorbox 列表样式 ==========
\tcbset{
    mylist/.style = {
        colframe = gray,
        colback = white,
        coltitle = red!50!yellow!3!white,
        colbacktitle = white,
        listing only,
        attach boxed title to top center = {yshift = -\tcboxedtitleheight/2},
        enhanced,
        drop fuzzy shadow,
        left = 6.5mm,
        breakable,
        fonttitle=\small\bfseries\color{black},
        before skip=12pt plus 2pt,
        after skip=12pt plus 2pt,
    },
    example/.style 2 args = {
        mylist,
        title = {Listing \thetcbcounter: #1},
        label = {#2},
    },
}

\newtcblisting[auto counter, number within = section, list inside = mcode]{matlab}[3][]{%
    listing options = {
        style = Matlab-editor,
        numbers = left,
        numberstyle = \footnotesize\color{darkgray}\ttfamily,
        basicstyle = \footnotesize\ttfamily,
        breaklines = true,
        breakatwhitespace = true,
        postbreak = \mbox{\textcolor{gray}{$\hookrightarrow$}\space}
    },
    overlay = {\begin{tcbclipinterior}\fill[blue!15!white] (frame.south west) rectangle ([xshift=5.3mm]frame.north west);\end{tcbclipinterior}},
    example = {#2}{#3}, #1,
}

\newtcblisting[auto counter, number within = section, list inside = pcode]{python}[3][]{%
    listing options = {
        language = Python,
        numbers = left,
        numberstyle = \footnotesize\color{darkgray}\ttfamily,
        basicstyle = \footnotesize\ttfamily,
        keywordstyle = \color{darkblue}\bfseries,
        commentstyle = \color{darkgreen},
        stringstyle = \color{red!60},
        showstringspaces = false,
        emphstyle = \color{orange!60},
        emph = {self, cls, True, False, None},
        breaklines = true,
        breakatwhitespace = true,
        postbreak = \mbox{\textcolor{gray}{$\hookrightarrow$}\space}
    },
    overlay = {\begin{tcbclipinterior}\fill[blue!15!white] (frame.south west) rectangle ([xshift=5.3mm]frame.north west);\end{tcbclipinterior}},
    example = {#2}{#3}, #1,
}

% ========== AI 报告专用无边框 listings 风格 ==========
\lstdefinestyle{ai_python_style}{
    language = Python,
    numbers = none,
    basicstyle = \footnotesize\ttfamily,
    keywordstyle = \color{darkblue}\bfseries,
    commentstyle = \color{darkgreen},
    stringstyle = \color{red!60},
    showstringspaces = false,
    emphstyle = \color{orange!60},
    emph = {self, cls, True, False, None},
    breaklines = true,
    breakatwhitespace = true,
    postbreak = \mbox{\textcolor{gray}{$\hookrightarrow$}\space},
    frame = none,
    backgroundcolor = \color{white},
    xleftmargin = 0pt,
    xrightmargin = 0pt,
}

\lstdefinestyle{ai_matlab_style}{
    style = Matlab-editor,
    numbers = none,
    basicstyle = \footnotesize\ttfamily,
    breaklines = true,
    breakatwhitespace = true,
    postbreak = \mbox{\textcolor{gray}{$\hookrightarrow$}\space},
    frame = none,
    backgroundcolor = \color{white},
    xleftmargin = 0pt,
    xrightmargin = 0pt,
}

\lstnewenvironment{ai_python}[1][]
    {\lstset{style=ai_python_style, #1}}
    {}

\lstnewenvironment{ai_matlab}[1][]
    {\lstset{style=ai_matlab_style, #1}}
    {}

% ========== 为定理类环境添加底色框 ==========
\mdfdefinestyle{thmshaded}{%
    backgroundcolor=gray!8,% 浅灰底色
    linecolor=black!30,% 边框颜色
    leftline=true,% 左侧线条
    topline=false,
    rightline=false,
    bottomline=false,
    linewidth=3pt,% 左侧线宽
    innerleftmargin=10pt,
    innertopmargin=6pt,
    innerbottommargin=6pt,
    roundcorner=0pt,
}

\surroundwithmdframed[style=thmshaded]{definition}
\surroundwithmdframed[style=thmshaded]{theorem}
\surroundwithmdframed[style=thmshaded]{lemma}
\surroundwithmdframed[style=thmshaded]{corollary}
\surroundwithmdframed[style=thmshaded]{assumption}
\surroundwithmdframed[style=thmshaded]{example}
\surroundwithmdframed[style=thmshaded]{problem}
\surroundwithmdframed[style=thmshaded]{solution}




\usepackage{tcolorbox} % 美丽盒子包
\usepackage{matlab-prettifier} % MATLAB代码排版包
\tcbuselibrary{listings,skins,breakable,xparse}

% 颜色定义
\definecolor{darkgray}{rgb}{0.4,0.4,0.4}
\definecolor{darkblue}{rgb}{0.0,0.0,0.5}
\definecolor{darkgreen}{rgb}{0,0.45,0.13}
\definecolor{orange}{rgb}{0.87,0.49,0.0}

% ========== tcolorbox 列表样式 ==========
\tcbset{
    mylist/.style = {
        colframe = gray,
        colback = white,
        coltitle = red!50!yellow!3!white,
        colbacktitle = white,
        listing only,
        attach boxed title to top center = {yshift = -\tcboxedtitleheight/2},
        enhanced,
        drop fuzzy shadow,
        left = 6.5mm,
        breakable,
        fonttitle=\small\bfseries\color{black},
        before skip=12pt plus 2pt,
        after skip=12pt plus 2pt,
    },
    example/.style 2 args = {
        mylist,
        title = {Listing \thetcbcounter: #1},
        label = {#2},
    },
}

\newtcblisting[auto counter, number within = section, list inside = mcode]{matlab}[3][]{%
    listing options = {
        style = Matlab-editor,
        numbers = left,
        numberstyle = \footnotesize\color{darkgray}\ttfamily,
        basicstyle = \footnotesize\ttfamily,
        breaklines = true,
        breakatwhitespace = true,
        postbreak = \mbox{\textcolor{gray}{$\hookrightarrow$}\space}
    },
    overlay = {\begin{tcbclipinterior}\fill[blue!15!white] (frame.south west) rectangle ([xshift=5.3mm]frame.north west);\end{tcbclipinterior}},
    example = {#2}{#3}, #1,
}

\newtcblisting[auto counter, number within = section, list inside = pcode]{python}[3][]{%
    listing options = {
        language = Python,
        numbers = left,
        numberstyle = \footnotesize\color{darkgray}\ttfamily,
        basicstyle = \footnotesize\ttfamily,
        keywordstyle = \color{darkblue}\bfseries,
        commentstyle = \color{darkgreen},
        stringstyle = \color{red!60},
        showstringspaces = false,
        emphstyle = \color{orange!60},
        emph = {self, cls, True, False, None},
        breaklines = true,
        breakatwhitespace = true,
        postbreak = \mbox{\textcolor{gray}{$\hookrightarrow$}\space}
    },
    overlay = {\begin{tcbclipinterior}\fill[blue!15!white] (frame.south west) rectangle ([xshift=5.3mm]frame.north west);\end{tcbclipinterior}},
    example = {#2}{#3}, #1,
}

% ========== AI 报告专用无边框 listings 风格 ==========
\lstdefinestyle{ai_python_style}{
    language = Python,
    numbers = none,
    basicstyle = \footnotesize\ttfamily,
    keywordstyle = \color{darkblue}\bfseries,
    commentstyle = \color{darkgreen},
    stringstyle = \color{red!60},
    showstringspaces = false,
    emphstyle = \color{orange!60},
    emph = {self, cls, True, False, None},
    breaklines = true,
    breakatwhitespace = true,
    postbreak = \mbox{\textcolor{gray}{$\hookrightarrow$}\space},
    frame = none,
    backgroundcolor = \color{white},
    xleftmargin = 0pt,
    xrightmargin = 0pt,
}

\lstdefinestyle{ai_matlab_style}{
    style = Matlab-editor,
    numbers = none,
    basicstyle = \footnotesize\ttfamily,
    breaklines = true,
    breakatwhitespace = true,
    postbreak = \mbox{\textcolor{gray}{$\hookrightarrow$}\space},
    frame = none,
    backgroundcolor = \color{white},
    xleftmargin = 0pt,
    xrightmargin = 0pt,
}

\lstnewenvironment{ai_python}[1][]
    {\lstset{style=ai_python_style, #1}}
    {}

\lstnewenvironment{ai_matlab}[1][]
    {\lstset{style=ai_matlab_style, #1}}
    {}

% ========== 为定理类环境添加底色框 ==========
\mdfdefinestyle{thmshaded}{%
    backgroundcolor=gray!8,% 浅灰底色
    linecolor=black!30,% 边框颜色
    leftline=true,% 左侧线条
    topline=false,
    rightline=false,
    bottomline=false,
    linewidth=3pt,% 左侧线宽
    innerleftmargin=10pt,
    innertopmargin=6pt,
    innerbottommargin=6pt,
    roundcorner=0pt,
}

\surroundwithmdframed[style=thmshaded]{definition}
\surroundwithmdframed[style=thmshaded]{theorem}
\surroundwithmdframed[style=thmshaded]{lemma}
\surroundwithmdframed[style=thmshaded]{corollary}
\surroundwithmdframed[style=thmshaded]{assumption}
\surroundwithmdframed[style=thmshaded]{example}
\surroundwithmdframed[style=thmshaded]{problem}
\surroundwithmdframed[style=thmshaded]{solution}


